\documentclass[12pt]{article}
% Préambule commun à tous les documents extraits de « La Revue CAPES »
\usepackage[T1]{fontenc}
\usepackage[utf8]{inputenc}
\usepackage{lmodern}
\usepackage{amsmath,amssymb,amsthm,amsfonts,mathrsfs}
\usepackage{setspace}
\usepackage{listings}
\usepackage{pgfplots}
\pgfplotsset{compat=1.18}
\usepackage{longtable}
\usepackage{adjustbox}
\usepackage{lettrine}
\usepackage{tkz-tab}
\usepackage{changepage}
\usepackage{pdflscape}
\usepackage{graphicx}
\usepackage{tikz}
\usepackage{xcolor}
\usepackage[a4paper,top=2cm,bottom=2.5cm,left=2.5cm,right=2cm]{geometry}
\usepackage{fancyhdr}
\usepackage{draftwatermark}
\SetWatermarkText{La RevueCapes}
\SetWatermarkLightness{0.9}
\SetWatermarkScale{2}
\usepackage{hyperref}
\hypersetup{colorlinks=true,linkcolor=blue!50!black,urlcolor=blue!50!black}

\definecolor{revue}{RGB}{20,60,120}

\pagestyle{fancy}
\fancyhf{}
\renewcommand{\headrulewidth}{0.4pt}
\setlength{\headheight}{15pt}

% \entete{Matière}{Titre}{Type de document}
\newcommand{\entete}[3]{%
  \fancyhead[L]{\small\textsc{La Revue CAPES} -- #1}%
  \fancyhead[R]{\small #2}%
  \fancyfoot[C]{\thepage}%
  \thispagestyle{plain}%
  \begin{center}
    {\color{revue}\rule{\textwidth}{1.2pt}}\\[4pt]
    {\small\textsc{Concours directs CAP-CEG / CAPES -- Burkina Faso}}\\[6pt]
    {\LARGE\bfseries\color{revue} #2}\\[6pt]
    {\large\itshape #1 -- #3}\\[2pt]
    {\color{revue}\rule{\textwidth}{1.2pt}}
  \end{center}
  \vspace{0.5cm}%
}

% Encadré pour les remarques ajoutées dans les corrigés
\newcommand{\remarque}[1]{\par\medskip\noindent\fcolorbox{revue}{revue!6}{\parbox{\dimexpr\linewidth-2\fboxsep-2\fboxrule}{\textbf{\color{revue}Remarque.} #1}}\par\medskip}


\begin{document}
\entete{Culture générale}{Corrigé du sujet type de dissertation n°14}{Correction}
\begin{spacing}{1.5}

\noindent\textbf{Sujet.} Quel rôle jouent le chômage des jeunes et l'absence d'éducation dans la propagation des groupes armés au Burkina Faso ?

\rubrique{1. Analyse du sujet}

\begin{itemize}
\item \textbf{Type de sujet :} question ouverte (« quel rôle ? ») qui demande d'expliquer un lien de cause à effet, puis d'en mesurer la portée (les autres facteurs) et de proposer des solutions.
\item \textbf{Mots-clés :}
  \begin{itemize}
  \item « chômage des jeunes » : absence d'emploi ou emploi précaire et insuffisant (sous-emploi), très répandu chez les jeunes ruraux comme chez les diplômés ;
  \item « absence d'éducation » : non-scolarisation, abandon scolaire, analphabétisme, mais aussi manque de formation professionnelle et d'éducation civique ;
  \item « propagation des groupes armés » : extension géographique des groupes terroristes et augmentation de leurs effectifs par le recrutement.
  \end{itemize}
\item \textbf{Problématique :} dans quelle mesure le chômage et le déficit d'éducation des jeunes favorisent-ils le recrutement et l'expansion des groupes armés, et comment rompre ce cercle vicieux ?
\end{itemize}

\rubrique{2. Plan détaillé}

\begin{enumerate}
\item[I.] \textbf{Le chômage des jeunes, un facteur de recrutement.} 1. La précarité économique et l'attrait des gains. 2. Le désœuvrement, la frustration et la quête de statut. 3. Les économies parallèles contrôlées par les groupes armés.
\item[II.] \textbf{L'absence d'éducation, un facteur de vulnérabilité.} 1. Le manque de compétences et de perspectives. 2. La vulnérabilité à l'endoctrinement. 3. Un cercle vicieux : l'insécurité détruit l'école.
\item[III.] \textbf{Des facteurs importants mais non exclusifs ; les solutions.} 1. D'autres causes (gouvernance, conflits, contexte régional, contrainte). 2. Créer des emplois et former les jeunes. 3. Rétablir l'école et éduquer à la paix.
\end{enumerate}

\rubrique{3. Proposition de rédaction}

\partie{Introduction}

Le Burkina Faso, comme les autres pays du Sahel, est confronté depuis plusieurs années à l'expansion de groupes armés terroristes qui recrutent principalement parmi les jeunes. Or, la jeunesse burkinabè, majoritaire dans la population, fait face à un chômage massif et à un accès souvent limité à l'éducation, surtout en milieu rural. On peut dès lors s'interroger sur le rôle que jouent le chômage des jeunes et l'absence d'éducation dans la propagation des groupes armés. Ces facteurs suffisent-ils à expliquer l'engagement de nombreux jeunes dans la violence ? Nous montrerons d'abord que le chômage constitue un puissant facteur de recrutement, puis que l'absence d'éducation rend les jeunes plus vulnérables, avant de souligner que ces facteurs, s'ils sont importants, ne sont pas les seuls, et de proposer des solutions.

\partie{I. Le chômage des jeunes, un facteur de recrutement}

\souspartie{1. La précarité économique et l'attrait des gains}
Dans de nombreuses zones rurales, les jeunes vivent d'une agriculture de subsistance soumise aux aléas climatiques, sans revenu régulier. Les groupes armés leur proposent de l'argent, une moto, un téléphone, une arme, parfois une prime d'engagement. Pour un jeune sans ressources, ces offres peuvent paraître attractives, d'autant qu'elles permettent de soutenir sa famille ou de se marier.

\souspartie{2. Le désœuvrement, la frustration et la quête de statut}
Le chômage n'est pas seulement une question d'argent : il prive le jeune de reconnaissance sociale et de perspectives d'avenir. Les groupes armés exploitent cette frustration en offrant un statut, un sentiment de pouvoir et d'appartenance, et en canalisant la colère contre l'État et les élites.

\souspartie{3. Les économies parallèles}
Les groupes armés contrôlent des activités lucratives : orpaillage, contrebande, trafic de carburant, vol de bétail. Ils y emploient des jeunes, qui deviennent progressivement dépendants de ces réseaux. La frontière entre économie de survie et engagement armé devient alors floue.

\transition{Le chômage explique en partie l'attrait des groupes armés. Mais il est souvent lié à un autre facteur : le déficit d'éducation.}

\partie{II. L'absence d'éducation, un facteur de vulnérabilité}

\souspartie{1. Le manque de compétences et de perspectives}
Un jeune qui n'a pas été à l'école ou qui l'a quittée tôt dispose de peu de compétences pour accéder à un emploi qualifié. Il reste enfermé dans la précarité, ce qui renforce l'effet du chômage. L'insuffisance de la formation professionnelle aggrave cette situation.

\souspartie{2. La vulnérabilité à l'endoctrinement}
L'éducation développe l'esprit critique et la connaissance de sa propre religion et de celle des autres. Un jeune peu instruit est plus facilement manipulé par des prédicateurs qui présentent une version déformée de la religion pour justifier la violence, ou par la propagande diffusée sur les réseaux sociaux.

\souspartie{3. Un cercle vicieux}
Les groupes armés ciblent eux-mêmes l'école, qu'ils considèrent comme un instrument de l'État : de nombreuses écoles ont été fermées, des enseignants menacés ou assassinés. Des centaines de milliers d'enfants sont ainsi privés de scolarité, ce qui prépare de futures recrues : l'insécurité détruit l'éducation, et l'absence d'éducation nourrit l'insécurité.

\transition{Chômage et absence d'éducation jouent donc un rôle important. Il faut cependant se garder d'en faire les seules causes du phénomène.}

\partie{III. Des facteurs importants mais non exclusifs ; quelles solutions ?}

\souspartie{1. D'autres causes interviennent}
Beaucoup de jeunes pauvres et peu instruits refusent la violence, et certains combattants sont instruits. D'autres facteurs comptent : l'absence de l'État dans certaines zones, les injustices et les abus, les conflits fonciers et communautaires, la vengeance après des exactions, le contexte régional et, souvent, la contrainte pure (enrôlements forcés, menaces contre les familles).

\souspartie{2. Créer des emplois et former les jeunes}
Il faut offrir aux jeunes des alternatives : formation professionnelle, appui à l'entrepreneuriat, financement de projets, valorisation de l'agriculture et de l'élevage, travaux à haute intensité de main-d'œuvre, encadrement de l'orpaillage. Un jeune qui a un avenir est moins tenté par la violence.

\souspartie{3. Rétablir l'école et éduquer à la paix}
Il faut rouvrir les écoles fermées dès que la sécurité le permet, développer l'éducation en situation d'urgence pour les élèves déplacés, et renforcer l'éducation civique, l'éducation à la tolérance et à la paix. Des programmes de désengagement et de réinsertion peuvent aussi aider les jeunes qui veulent quitter les groupes armés.

\partie{Conclusion}

Le chômage des jeunes et l'absence d'éducation jouent un rôle important dans la propagation des groupes armés au Burkina Faso : le premier pousse des jeunes sans ressources vers des groupes qui leur offrent revenus et statut, la seconde les rend plus vulnérables à la manipulation et prive le pays de perspectives d'avenir. Ces facteurs ne sont cependant pas les seuls : la faiblesse de l'État, les injustices, les conflits locaux et la contrainte y contribuent également. Pour tarir le recrutement des groupes armés, il faut donc compléter l'action militaire par une politique ambitieuse en faveur de l'emploi et de l'éducation des jeunes. La jeunesse ne doit pas être considérée comme une menace, mais comme la principale ressource pour construire la paix.

\end{spacing}
\end{document}
